\newpage
\appendix
\section{Anexo 1: Proceso de Generación de Ideas y Selección del Proyecto}

La definición de la temática definitiva del presente Trabajo Profesional fue el resultado de una dinámica de deconstrucción y selección basada en un ejercicio de \textit{Brainwriting} intrateam. Mediante esta técnica de generación de ideas silenciosa, iterativa y cruzada, el equipo planteó diversas soluciones tecnológicas orientadas a resolver fallas de mercado o vacíos operativos específicos. Tras una primera fase de filtrado por consenso, se preseleccionaron cuatro propuestas principales:

\begin{enumerate}
    \item \textbf{ParkLink:} Plataforma web/móvil para el mapeo, gestión y reserva manual de cocheras urbanas ociosas conectando conductores y propietarios particulares.
    \item \textbf{Prompt Coaching:} Middleware analítico orientado a optimizar la interacción humano-IA, reduciendo el consumo de tokens y detectando puntos débiles en los prompts.
    \item \textbf{Consorcio360:} Plataforma digital SaaS para la gestión administrativa de consorcios con análisis de gastos mediante IA y un marketplace integrado de profesionales.
    \item \textbf{VoxLedger:} Chatbot conversacional para el registro y estructuración de flujos financieros, deudas y stock mediante mensajes de voz, enfocado en la economía informal.
\end{enumerate}

\subsection{Análisis de Viabilidad: Las 5Cs de las Alternativas}
A fin de evaluar estratégicamente cada alternativa bajo un marco unificado, se aplicó el modelo de las 5Cs (\textit{Company, Customers, Competitors, Collaborators, Context}):

\begin{itemize}
    \item \textbf{ParkLink:}
    \begin{itemize}
        \item \textit{Company:} Aplicación y web app para reservar, alquilar y gestionar espacios públicos o privados disponibles para estacionamiento por la ciudad.
        \item \textit{Customers:} Conductores particulares urbanos y propietarios de plazas de estacionamiento subutilizadas.
        \item \textit{Competitors:} Blinkay (estacionamiento medido municipal),
        parkings tradicionales, MercadoLibre y OLX (alquiler mensual de cocheras),
        Waze y Google Maps (información de tráfico sin reserva), y la búsqueda
        errática en calle.
        \item \textit{Collaborators:} Consorcios residenciales, administraciones locales, tráfico urbano, mal manejo del sistema de estacionamiento medido.
        \item \textit{Context:} Crecimiento exponencial del parque automotor en CABA, saturación de la vía pública y digitalización de micro-pagos.
    \end{itemize}
    \item \textbf{Prompt Coaching:}
    \begin{itemize}
        \item \textit{Company:} Middleware que analiza la interacción entre un usuario y un agente de inteligencia artificial, y plantea mejoras para que el usuario pueda identificar cuáles son sus puntos débiles.
        \item \textit{Customers:} Industrias de software (ej. Mercado Libre, Accenture), estudios jurídicos y agencias de marketing masivo.
        \item \textit{Competitors:} Generadores de prompts integrados, cursos de ingeniería de prompts y funciones nativas de optimización de los propios proveedores de LLMs.
        \item \textit{Collaborators:} Interfaces de desarrollo (VS Code, Cursor) y proveedores de modelos (OpenAI, Anthropic).
        \item \textit{Context:} Incertidumbre regulatoria global a corto plazo sobre la retención de datos conversacionales corporativos y la propiedad del output de IA.
    \end{itemize}
    \item \textbf{Consorcio360:}
    \begin{itemize}
        \item \textit{Company:} Plataforma digital SaaS (Software as a Service) orientada a la gestión de consorcios. Su objetivo es modernizar la administración tradicional mediante la digitalización de procesos y la incorporación de herramientas de análisis inteligente.
        \item \textit{Customers:} Administradores de consorcios tradicionales y prestadores de servicios de mantenimiento.
        \item \textit{Competitors:} Plataformas establecidas (ConsorcioAbierto, Octopus Proptech), planillas de Excel y redes informales de contacto.
        \item \textit{Collaborators:} Proveedores de infraestructura en la nube, procesadores de pago y profesionales matriculados (electricistas, plomeros).
        \item \textit{Context:} Alta inflación y necesidad urgente de los consorcios por auditar sobrecostos, contrarrestada por una fuerte resistencia cultural y burocrática al cambio tecnológico.
    \end{itemize}
    \item \textbf{VoxLedger:}
    \begin{itemize}
        \item \textit{Company:} Aplicación y chatbot (integrable con WhatsApp, Telegram o web app) que permite a los usuarios registrar pedidos, gastos, ingresos y stock mediante audios o mensajes de texto.
        \item \textit{Customers:} Emprendedores independientes, freelancers, trabajadores de oficios y comercios barriales informales.
        \item \textit{Competitors:} Anotaciones manuales en papel, notas de voz de WhatsApp sin estructurar y hojas de cálculo tradicionales.
        \item \textit{Collaborators:} Proveedores de modelos Speech-to-Text, redes de mensajería (WhatsApp, Telegram) y ecosistemas fintech.
        \item \textit{Context:} Elevada tasa de informalidad económica regional y precarización digital que demanda interfaces de fricción nula (conversacionales).
    \end{itemize}
\end{itemize}
Tras este análisis preliminar, la propuesta de \textbf{Prompt Coaching} fue descartada debido a su alta dependencia de la evolución de terceros (los proveedores de IA reducen la necesidad de optimizadores externos a medida que sus modelos interpretan mejor la intención natural) y las severas restricciones del contexto legal corporativo.

\vspace{0.5cm}
\begin{table}[H]
\centering
\renewcommand{\arraystretch}{2.2}
\resizebox{\textwidth}{!}{%
\begin{tabular}{|l|c|c|c|c|c|}
\hline
\textbf{Proyecto} & \textbf{Company} & \textbf{Customers} & \textbf{Competitors} & \textbf{Collaborators} & \textbf{Context} \\
\hline
\textbf{ParkLink}         & \gcircle & \gcircle & \gcircle & \ycircle & \gcircle \\
\hline
\textbf{Prompt Coaching}  & \ycircle & \ycircle & \rcircle & \ycircle & \ycircle \\
\hline
\textbf{Consorcio360}     & \ycircle & \ycircle & \rcircle & \gcircle & \ycircle \\
\hline
\textbf{VoxLedger}        & \gcircle & \gcircle & \gcircle & \ycircle & \gcircle \\
\hline
\end{tabular}%
}
\caption{Análisis comparativo de las 5Cs por propuesta}
\label{tab:5cs}
\end{table}

\noindent\gcircle~Favorable \quad \ycircle~Moderado \quad \rcircle~Desfavorable

\subsection{Estrategia de Mercado y Proyecciones: Las 4Ps de los Productos Finalistas}
Para los tres productos que superaron el filtro inicial, se modeló su estrategia de comercialización (\textit{Product, Price, Place, Promotion}) junto con una estimación cuantitativa de adopción para el primer año en el mercado argentino.



\subsubsection{Consorcio360}
\begin{itemize}
    \item \textbf{Product:} Plataforma web y móvil SaaS con dashboard financiero, liquidación de expensas, reserva de espacios comunes, marketplace de proveedores calificados y auditoría de sobrecostos mediante IA.
    \item \textbf{Price:} Esquema híbrido. Plan \textit{Freemium} con gestión básica. Plan \textit{Premium} escalonado según unidades funcionales (Ticket promedio estimado de \$40.000 ARS/mes por consorcio) más una comisión del 5\% por contrataciones en el marketplace.
    \item \textbf{Place:} Despliegue en CABA como mercado inicial, expandiéndose hacia el Gran Buenos Aires y los principales centros urbanos del país.
    \item \textbf{Promotion:} Venta directa a administradores mediante LinkedIn, pauta digital segmentada en Google Ads e informes mensuales de tendencias de costos por zona para generar autoridad de marca.
    \item \textbf{Proyección Financiera (Año 1):} Con un universo nacional de 120.000 consorcios, el escenario medio estima alcanzar 600 consorcios activos, proyectando un ingreso combinado anual de \$282.857 USD (\$205.714 USD vía suscripción SaaS y \$77.143 USD por comisiones del marketplace).
\end{itemize}

\subsubsection{VoxLedger}
\begin{itemize}
    \item \textbf{Product:} Chatbot inteligente integrado a canales de mensajería instantánea que procesa notas de voz informales, extrae datos operativos de manera automática y alimenta un panel de métricas de caja e inventario sin carga manual de formularios.
    \item \textbf{Price:} Modelo \textit{Freemium}. Plan gratuito limitado en volumen de audios mensuales. Plan \textit{Premium} de \$8.000 ARS mensuales con almacenamiento histórico completo, múltiples dispositivos vinculados y reportes automáticos exportables.
    \item \textbf{Place:} Distribución digital global focalizada inicialmente en mercados hispanohablantes con altos índices de informalidad laboral (Argentina, México, Colombia y Perú).
    \item \textbf{Promotion:} Demostraciones dinámicas de flujo punta a punta (audio a gráfico) en formatos de video corto (YouTube Shorts, TikTok) y campañas orgánicas orientadas a comunidades de emprendedores.
    \item \textbf{Proyección Financiera (Año 1):} Apuntando a una conversión conservadora del 8\% sobre una base estimada de 20.000 usuarios activos en el escenario medio (1.600 usuarios Premium), se proyecta un ingreso anual de \$109.714 USD.
\end{itemize}

\subsubsection{ParkLink}
\begin{itemize}
    \item \textbf{Product:} Ecosistema digital (Web/Mobile) dedicado al mapeo cartográfico georreferenciado y a la reserva programada y manual de espacios de estacionamiento privados infrautilizados.
    \item \textbf{Price:} Modelo transaccional puro: comisión del 10\% sobre el valor de cada reserva descontada al propietario del espacio, más un cargo de servicio del 5\% aplicado al conductor. Publicar un espacio es gratuito; ParkLink solo genera ingresos cuando hay una transacción exitosa.
    \item \textbf{Place:} Concentración absoluta inicial en los corredores de alta densidad vehicular y comercial de la Ciudad Autónoma de Buenos Aires (CABA).
    \item \textbf{Promotion:} Campañas digitales geocercadas dirigidas a conductores recurrentes, programas de referidos vecinales e incentivos para la captación temprana de oferentes residenciales.
    \item \textbf{Proyección Financiera (Año 1):} Estimando un valor de reserva promedio de \$6.000 ARS por estadía corta (comisión neta de \$600 ARS), el escenario intermedio (1.0\% de penetración de mercado con 2.160 reservas diarias) proyecta una recaudación anual para la plataforma de \$333.257 USD.
\end{itemize}

\vspace{0.5cm}
\begin{table}[H]
\centering
\renewcommand{\arraystretch}{2.2}
% Reducir el tamaño de fuente para igualar la tabla de las 5Cs
{\small
\resizebox{0.85\textwidth}{!}{%
\begin{tabular}{|l|c|c|c|c|}
\hline
\textbf{Proyecto} & \textbf{Product} & \textbf{Price} & \textbf{Place} & \textbf{Promotion} \\
\hline
\textbf{Consorcio360} & \ycircle & \ycircle & \ycircle & \ycircle \\
\hline
\textbf{VoxLedger}     & \ycircle & \ycircle & \gcircle & \gcircle \\
\hline
\textbf{ParkLink}      & \gcircle & \gcircle & \gcircle & \gcircle \\
\hline
\end{tabular}%
}
}
\caption{Viabilidad comparativa de las 4Ps por propuesta}
\label{tab:4ps}
\end{table}

\noindent\gcircle~Favorable \quad \ycircle~Moderado \quad \rcircle~Desfavorable


\subsection{Justificación de la Selección de ParkLink}
La selección unánime de \textbf{ParkLink} sobre \textit{Consorcio360} y \textit{VoxLedger} como eje de este Trabajo Profesional responde a tres decisiones de diseño de ingeniería y control de riesgos del proyecto:

\begin{enumerate}
    \item \textbf{Autonomía del MVP e Integración Abierta:} A diferencia de \textit{Consorcio360}, cuya viabilidad comercial e implementación requiere romper barreras culturales de administración sumamente rígidas y procesos de toma de decisión lentos en asambleas de copropietarios, ParkLink funciona bajo un modelo P2P directo de baja fricción y adopción inmediata. Asimismo, frente a \textit{VoxLedger}, evita la dependencia absoluta de canales de mensajería de terceros (como las políticas de las APIs corporativas de Meta para WhatsApp).
    \item \textbf{Control de la Complejidad y Viabilidad del Alcance Académico:} Al acotar estrictamente el alcance de ParkLink al mapeo interactivo y la gestión de reservas manuales programadas (postergando deliberadamente la automatización mediante sensores físicos de ocupación o algoritmos predictivos complejos para fases futuras del ciclo de vida del producto), se asegura un plan de desarrollo realizable por el equipo dentro de los tiempos del ciclo lectivo, garantizando robustez de software por sobre la dispersión tecnológica.
    \item \textbf{Previsibilidad y Validación de Datos Financieros:} El modelo transaccional de ParkLink ofrece una monetización lineal directamente ligada a la resolución del problema físico del conductor (tiempo de búsqueda de estacionamiento). Esto permite una validación de métricas de negocio más transparente y medible dentro de un entorno urbano acotado como CABA, en comparación con la incertidumbre en las tasas de conversión de los modelos \textit{freemium} proyectados en las otras alternativas.
\end{enumerate}
